%%%%%     PACKS     %%%%%
\documentclass[12pt]{article}
\usepackage[margin=1in,headsep=.60in]{geometry}
\usepackage[utf8]{inputenc}
\usepackage[table, dvipsnames]{xcolor}
\usepackage{array}
\usepackage{amsmath}
\usepackage{booktabs}
\usepackage{mdframed} %For box around text
\usepackage{hyperref} %For hyperlinks in the PDF
\hypersetup{
    colorlinks=true,
    linkcolor=blue,
    filecolor=magenta,
    urlcolor=cyan,
}
\urlstyle{same}
\usepackage{amsthm}
\usepackage{amssymb}
\usepackage{amsfonts}
\usepackage{siunitx}
\usepackage{graphicx}
\usepackage{caption}
\usepackage{pgfplots}
\graphicspath{{Images/}}
\usepackage[colorinlistoftodos]{todonotes}
\usepackage{cleveref}
\usepackage[labelformat=simple]{subcaption}
\usepackage{grffile}
\usepackage{gensymb}
\usepackage{float}
\usepackage[shortlabels]{enumitem}
\usepackage{enumitem}
\setlistdepth{9}
\usepackage{biblatex} %Imports biblatex package
\addbibresource{references.bib}
\usepackage{outlines}
\usepackage{minted}
\usepackage{pdflscape}
\usepackage{everypage}
\usepackage{multirow}
\usepackage{multicol}
\usepackage{afterpage}
\usepackage[labelfont=bf, font=small]{caption}
\usepackage{upgreek}
\usepackage{tikz}
\usepackage{xcolor}
\usepackage{listings}

% style sheet for code
\definecolor{codegreen}{rgb}{0,0.6,0}
\definecolor{codegray}{rgb}{0.5,0.5,0.5}
\definecolor{codepurple}{rgb}{0.58,0,0.82}
\definecolor{backcolour}{rgb}{0.95,0.95,0.92}

\lstdefinestyle{mystyle}{
    backgroundcolor=\color{backcolour},   
    commentstyle=\color{codegreen},
    keywordstyle=\color{magenta},
    numberstyle=\tiny\color{codegray},
    stringstyle=\color{codepurple},
    basicstyle=\ttfamily\footnotesize,
    breakatwhitespace=false,         
    breaklines=true,                 
    captionpos=b,                    
    keepspaces=true,                 
    numbers=left,                    
    numbersep=5pt,                  
    showspaces=false,                
    showstringspaces=false,
    showtabs=false,                  
    tabsize=2
}

\lstset{style=mystyle}

\usepackage{lastpage}
%%%%%     COMMANDS     %%%%%
%%%% Blue box for subsection text (not figures) %%%%
\newenvironment{bluebox}
  {\begin{mdframed}[backgroundcolor=blue!5,linecolor=blue!40,roundcorner=8pt]}
  {\end{mdframed}}

%%%% Simple neutral box for figures %%%%
\newenvironment{figbox}
  {\begin{mdframed}[roundcorner=8pt,shadow=true,shadowsize=4pt,shadowcolor=black!40]}
  {\end{mdframed}}

\begin{document}

\begin{titlepage}
\newcommand{\HRule}{\rule{\linewidth}{0.5mm}}

\center

\textsc{\LARGE Aarhus University}\\[1.5cm]
\textsc{\Large Department Of Computer Science }\\[0.5cm]
\textsc{\large Compilation}\\[0.5cm]
    

\HRule\\[0.4cm]
    \center 
    {\huge\bfseries  Hand-in 3 (Individual Report)}\\[0.4cm] % Title of your document
\HRule\\[1.5cm]

\begin{minipage}{0.4\textwidth}
        \begin{flushleft}
            \large
            \textit{Author}\\   
            Søren M. \textsc{Damsgaard}\\
                 % Your name
        \end{flushleft}
    \end{minipage}
~
    \begin{minipage}{0.4\textwidth}
        \begin{flushright}
            \large
            \textit{Student number}\\
            \textbf{202309814}\\
                 % Studienummer\
            
        \end{flushright}
    \end{minipage}

\vfill\vfill\vfill % Position the date 3/4 down the remaining page
    
    {\large\today}

\vfill\vfill
    \includegraphics[width=0.2\textwidth]{Aarhus_University_seal.png}\\[1cm] 
    

\vfill
\end{titlepage}
%%%%%      CHAPTERS      %%%%%
\hspace{0.02cm}
\section*{Task 1 - Implementation of Semantic Analysis}
The implementation of this heavily focuses on the module \texttt{semant.ml} and it's functions \texttt{infertype\_expr} and \texttt{typecheck\_expr}
, these functions recursively call each other to traverse the AST translating the untyped AST into a typed AST, making sure to follow Dolphins type rules.\\
This implementation also uses the modules \texttt{env} and \texttt{errors} to handle variable scopes when encountering nested variables and handling errors/mismatched types.
\section*{Task 2 - Implementation of Code Generation}
Generating code by translating the TypedAst to functional LLVM-\- code. is a huge part of this project, hence a seperate module \texttt{codegen.ml} was created\\
This was doine by using the CfgBuilder to create basic blocks, and allocating variables using \texttt{Ll.Alloca} with a the cgenv environment, that maps symbols to LLVM pointers.\\
In the code there is an explicit step when using a division or remainder op, it's a \texttt{Icmp Eq} check that safely exits if a divisorr is zero.
\section*{Task 3 - Testing }
The module \texttt{compiler.ml} takes an untyped AST program, sends it to the semantic analysis, and if the error list is empty it continues
to the codegen module that returns a LLVM-\- program, and if an error is found the list of errors is printed and the compile is aborted.\\
To be extra safe, an entire test corpus was made, with A LOT of tests that covers everything, both trivial and more concrete features.

\section*{Questions}

\subsection*{Question 1 - type Inference or type checking }
Type inference is used when a variable declaration doesn't have an explicit type, and it has to be infere from the bottom-up in the AST.\\
Type checking is when the semantics has strict rules, such as arithmetic operators (Plus, Minus) that needs the l/r operands to be Integers.\\
 Combined, inference and checking supports eachother well.

\subsection*{Question 2}
The way we implementede error recovery was to make a dedicated ErrorType  if an AST node had a mismatched type.\\
When the \texttt{typecheck\_expr} encounters a mismatch it appends the error to an Error list and returns the expression with a \texttt{TAst.ErrorType} continuing the analysis,
 and in the end it logs all errors/mismatch, making it easier to see what went wrong.
 
\subsection*{Question 3}
*Glory - Didn't do it.
\section*{How to run the code}
You will have received a zip file containing the complete compiler source code, for phase 1, the Dune build configurations, and the C runtime file.\\
\newline
This is run/compiled in VScode using the docker container provided.\\
To compile and run the project inside a VS Code development container, execute the following commands in the terminal while inside the folder:
\begin{enumerate}
    \item \texttt{dune build} 
    \item \texttt{dolphin compile --from-ast --phase 1 tests/<test name>.json -o out.ll} (compiles the serialized AST and outputs an LLVM- - file)
    \item \texttt{clang runtime.c out.ll -o res} 
    \item \texttt{./res} (runs the progran and outputs any intermediat prints in the terminal)
    \item \texttt{echo \$?} (outputs the return statement from he LLVM-- code to the terminal)
\end{enumerate}
\end{document}